\documentclass[11pt,letterpaper]{article}
\usepackage[utf8]{inputenc}
\usepackage[T1]{fontenc}
\usepackage{amsmath,amsfonts,amssymb}
\usepackage{bm}
\usepackage{graphicx}
\usepackage{geometry}
\usepackage{fancyhdr}
\usepackage{hyperref}
\usepackage{booktabs}
\usepackage{array}
\usepackage{float}
\usepackage{caption}
\usepackage{xcolor}

\geometry{margin=1in}
\setlength{\headheight}{14pt}
\pagestyle{fancy}
\fancyhf{}
\rhead{Contact stability of a monolithic Li-ion cell}
\lhead{Preliminary chemo-mechanical study}
\cfoot{\thepage}

\hypersetup{colorlinks=true,linkcolor=blue,urlcolor=blue,citecolor=blue}

\title{\textbf{Mechanical contact-stability of a monolithic, pressure-free
Li-ion cell: a preliminary chemo-mechanical study}}
\author{}
\date{\today}

\begin{document}
\maketitle

\begin{abstract}
We present a preliminary, mechanics-only continuum model of a monolithic
Li-ion cell engineered to operate without external stack pressure.  Two
polymerization-driven volumetric eigenstrains---one in the gel polymer
electrolyte and one in the conformal package coating---together with
optional external stack pressure and lithiation-induced electrode
breathing, generate the through-thickness contact pressure required for
low-resistance interfaces.  We close the model with an exponential
pressure--resistance relation $R_c(p_c) = R_0 \exp(-\alpha p_c)$ and
report a contact-stability metric $\Pi_c = \langle p_c \rangle /
p_{\min}$, with $p_{\min}$ set so that $R_c(p_{\min})$ equals the cell
target.  A four-axis design grid (GPE shrink, package shrink, applied
stack pressure, GPE modulus) reveals a clearly bounded feasible window
where $\Pi_c \ge 1$ is achievable.  Package shrink is the dominant
linear lever; GPE shrink mildly unloads the separator interfaces and
must be co-designed with package shrink to avoid local debonding;
stiffening the GPE provides a diminishing-returns boost that
saturates near $E_{\rm sep} \approx 25$~MPa.  Lithiation breathing
modulates $\Pi_c$ by $\sim$25\% per cycle, so the design must clear
target at the discharged state.
\end{abstract}

\section{Scope}

The cell concept of interest is a monolithic Li-ion stack in which a
continuous polymer network binds the porous current collectors,
composite electrodes, and separator into a single body, with a
conformal polymer + ceramic coating providing the hermetic package and
the source of mechanical preload.  The design objective is to maintain
adequate through-thickness contact pressure---and thus low contact
resistance---without applying external stack pressure during cycling.

This study addresses one specific question: \emph{can the
polymerization-induced eigenstrains in the GPE and the package
envelope, together with optional modest external pressure, deliver
enough average contact pressure to land the contact resistance below
target?}  We deliberately restrict to mechanics and a closed-form
$R_c(p_c)$ relation so that the central output, the contact-stability
metric $\Pi_c$, can be evaluated cheaply across a four-dimensional
design grid.  Coupled ionic and electronic transport are deferred to a
follow-on study.

\section{Geometry}

We model a 2D plane-strain cross-section of the pouch.  Bottom-to-top
the active stack is
\[
\text{Cu collector}\;|\;\text{anode}\;|\;\text{separator (GPE)}\;|\;\text{cathode}\;|\;\text{Al collector},
\]
fully enveloped on all four sides by a polymer package layer.
Per-stack layer thicknesses (Table~\ref{tab:geom}) follow a
representative phase-3 cell-design specification.  Total package
thickness is $t_{\rm pkg} = 100~\mu$m on each face.

The wrapped envelope is what generates through-thickness compression in
this model.  With lateral rollers (plane-strain mirror symmetry), the
\emph{side} panels of the package envelope contract along $y$ as they
cure; because their bottom and top edges connect to the bottom and top
package strips, this contraction pulls those strips inward and
compresses the active stack.  A geometry with only top/bottom strips
and a free top boundary gives $\sigma_{yy}\equiv 0$ throughout (force
balance on the free surface), so the closed wrap is an essential
ingredient of the mechanism.

\paragraph{Caveat on slice width.}
The slice width is set to $W_{\rm inner} = 1$~mm in this preliminary
study---a numerical-economy choice.  For the realistic cell footprint
($\sim$$4 \times 4$~cm$^2$) the actual lateral extent is
$W_{\rm inner} \approx 40$~mm, and the side-panel wrap mechanism scales
as $t_{\rm pkg} / W_{\rm inner}$.  At $W_{\rm inner} = 1$~mm this ratio
is $0.10$; at the realistic value it is $0.0025$ (i.e.\ $\sim$40$\times$
smaller).  Plane-strain modeling also captures only one of the two
in-plane wrap directions, contributing another factor of $\sim$2.  The
absolute $\Pi_c$ values reported below should therefore be read as
\emph{relative} indicators along the design axes.  A 3D shell-wrap
model, or a 2D analysis with $W_{\rm inner}$ set to the actual lateral
extent, would be required before the absolute numbers can be quoted.

\begin{figure}[H]
\centering
\includegraphics[width=0.88\linewidth]{figures/geometry_schematic.pdf}
\caption{Cross-section of the wrapped pouch.  Hatched region is the
polymer package envelope (top, bottom, and both side panels);
non-hatched layers are the active stack.  Red arrows indicate the two
cure-shrinkage eigenstrains; blue arrows mark the symmetry-roller and
fixed-bottom boundary conditions; green arrows mark the optional
external stack pressure $\sigma_{\rm stack}$ on the top face.}
\label{fig:geom}
\end{figure}

\begin{table}[H]
\centering
\caption{Per-stack layer thicknesses [mm].}
\label{tab:geom}
\begin{tabular}{lcc}
\toprule
Layer & Symbol & Default \\
\midrule
Cu current collector (porous) & $t_{\rm cu}$       & 0.025 \\
Graphite anode                & $t_{\rm anode}$    & 0.055 \\
GPE / separator               & $t_{\rm sep}$      & 0.055 \\
LVPF cathode                  & $t_{\rm cathode}$  & 0.137 \\
Al current collector (porous) & $t_{\rm al}$       & 0.025 \\
Package layer (each face)     & $t_{\rm pkg}$      & 0.100 \\
\bottomrule
\end{tabular}
\end{table}

\section{Constitutive model}

Each block is treated as a homogenized, isotropic, compressible
neo-Hookean solid.  We adopt a multiplicative split of the deformation
gradient
\begin{equation}
\mathbf{F} = \mathbf{F}_m \, \mathbf{F}_g, \qquad
\mathbf{F}_g = \mathbf{F}_c \, \mathbf{F}_l ,
\label{eq:multi-split}
\end{equation}
where $\mathbf{F}_c$ is the cure-shrinkage eigen part, $\mathbf{F}_l$
is the lithiation eigen part, and $\mathbf{F}_m$ carries the elastic
response.  Only $\mathbf{F}_m$ enters the energy density,
\begin{equation}
\psi_m(\mathbf{F}_m) = \tfrac{\lambda}{2}(\ln J_m)^2 - G \ln J_m
+ \tfrac{G}{2}(\mathbf{F}_m\!:\!\mathbf{F}_m - 3),
\quad J_m = \det \mathbf{F}_m,
\end{equation}
with first Piola--Kirchhoff stress $\mathbf{P} = \partial
\psi_m/\partial \mathbf{F}_m \cdot \mathbf{F}_g^{-T}$ and Cauchy stress
$\bm{\sigma} = J^{-1}\,\mathbf{P}\,\mathbf{F}^T$.  Eigenstrains and
elastic strains are both small ($\lesssim$~10\%), so the geometric
corrections are modest and the analysis is performed on the undeformed
mesh.

Default Lam\'e parameters per block are derived from $E$ and $\nu$ in
Table~\ref{tab:moduli}.

\begin{table}[H]
\centering
\caption{Default linear-elastic moduli per block.  Values are
representative; the GPE modulus is varied as a fourth design axis (see
Section~\ref{sec:design}).}
\label{tab:moduli}
\begin{tabular}{lccl}
\toprule
Block          & $E$ [MPa] & $\nu$ & Notes \\
\midrule
Cu collector   & $5.0\times 10^{4}$ & 0.34 & Porous-homogenized \\
Graphite anode & $1.0\times 10^{3}$ & 0.30 & Composite electrode \\
GPE/separator  & 10                 & 0.45 & Soft, near-incompressible \\
LVPF cathode   & $5.0\times 10^{3}$ & 0.30 & Composite electrode \\
Al collector   & $3.0\times 10^{4}$ & 0.33 & Porous-homogenized \\
Package layer  & $1.0\times 10^{3}$ & 0.40 & Acrylate + ceramic \\
\bottomrule
\end{tabular}
\end{table}

\section{Eigenstrain physics}

\subsection{Cure-induced shrinkage}

Polymerization of the in-situ-cured GPE and of the conformal package
coating is accompanied by a volumetric shrinkage of 2--14\%, typical
of in-situ-polymerized acrylates.  We model both as an isotropic,
cure-driven eigen deformation gradient,
\begin{equation}
\mathbf{F}_c = \sqrt[3]{1 - c\,\varepsilon_{\rm sh}}\;\mathbf{I},
\qquad c \in [0,1],
\end{equation}
where $c$ is a dimensionless cure variable that ramps from 0 to 1 over
a pseudo-time interval representing the polymerization step.  The
volumetric shrinkage $\varepsilon_{\rm sh}$ is block-dependent:
$\varepsilon_{\rm GPE}$ in the separator, $\varepsilon_{\rm pkg}$ in
the package, and zero in the conductor and electrode blocks.

\subsection{Lithiation-induced breathing}

During cycling, the active electrodes breathe.  We model this as a
state-of-charge-driven eigen deformation gradient,
\begin{equation}
\mathbf{F}_l = \sqrt[3]{1 + \alpha_{\rm Li}\,s}\;\mathbf{I},
\qquad s \in [0,1] ,
\end{equation}
with sign convention: $\alpha_{\rm Li} > 0$ for materials that expand
on lithiation, $\alpha_{\rm Li} < 0$ for materials that contract on
lithiation.  $s = 0$ is fully discharged (cathode lithiated, anode
delithiated) and $s = 1$ is fully charged.  Defaults are
$\alpha_{\rm Li}^{\rm anode} = +0.10$ (graphite expands $\sim$10\% on
full lithiation) and $\alpha_{\rm Li}^{\rm cathode} = -0.05$
(low-volume-change LVPF).  The two eigen parts compose multiplicatively
into the total eigen deformation $\mathbf{F}_g = \mathbf{F}_c
\mathbf{F}_l$ in Eq.~\eqref{eq:multi-split}.

Figure~\ref{fig:cycling} shows the contact-stability response under a
two-cycle charge/discharge sweep at the baseline shrink combination
($\varepsilon_{\rm GPE} = 0.05$, $\varepsilon_{\rm pkg} = 0.05$,
$\sigma_{\rm stack} = 0$).  After cure, $\Pi_c$ oscillates between a
post-cure floor at the discharged state ($s = 0$, $\Pi_c\approx 0.68$)
and a peak $\sim$25\% higher at full charge ($s = 1$,
$\Pi_c\approx 0.84$): the anode swells more than the cathode shrinks,
the stack thickens against the rigid envelope, and contact pressure
rises.  At this baseline shrink combination the cell sits below the
$\Pi_c = 1$ target throughout the cycle; engineering toward higher
$\varepsilon_{\rm pkg}$ or modest $\sigma_{\rm stack}$ lifts both the
floor and the peak across unity.

\begin{figure}[H]
\centering
\includegraphics[width=0.92\linewidth]{figures/cycling_trajectory.pdf}
\caption{Cycling response at the baseline design point.  Top panel:
prescribed state of charge $s$.  Bottom panel: contact-stability
metric $\Pi_c$ (left axis, navy) and stack-averaged contact pressure
$\langle p_c \rangle$ (right axis, red dashed).  Shaded region is the
cure ramp.}
\label{fig:cycling}
\end{figure}

\section{Boundary conditions and loading}

\begin{itemize}
\item \textbf{Lateral edges} ($x = 0,\,W_{\rm outer}$): roller,
$u_x = 0$.  Models in-plane mirror symmetry of a long pouch and
prevents rigid-body translation; with $u_x=0$ on both sides the package
side panels' volumetric shrinkage is fully redirected into the $y$
direction.
\item \textbf{Bottom edge} ($y = 0$): roller, $u_y = 0$.
\item \textbf{Top edge} ($y = H_{\rm outer}$): traction
$t_y = -\sigma_{\rm stack}$, where $\sigma_{\rm stack} \ge 0$ is the
externally applied stack pressure.  $\sigma_{\rm stack} = 0$ is the
pressure-free baseline.
\end{itemize}

Cure is ramped from 0 to 1 over $t \in [0, 1]$, then plateaus.  After
cure completes, $s$ may cycle sinusoidally,
\begin{equation}
s(t) = \tfrac{1}{2}A_{\rm cyc} \bigl[1 - \cos\!\bigl(2\pi(t-1)/T_{\rm cyc}\bigr)\bigr],
\quad t \ge 1,
\end{equation}
with amplitude $A_{\rm cyc}$ and period $T_{\rm cyc}$, sufficient to
demonstrate that contact pressure tracks lithiation and to bound the
worst-case $\Pi_c$ over a cycle.

\section{Contact-resistance closure and stability metric}

For each electrode interface $\Gamma_k$ we define an average normal
contact pressure
\begin{equation}
p_c^{(k)} = - \frac{1}{|\Gamma_k|} \int_{\Gamma_k} \sigma_{yy}\,\mathrm{d}A
\end{equation}
(positive in compression).  The contact resistance per interface
follows the standard exponential closure
\begin{equation}
R_c(p_c) = R_0 \, \exp\!\bigl(-\alpha\, p_c\bigr),
\label{eq:Rc}
\end{equation}
with $R_0 = 200$~m$\Omega\cdot$cm$^2$ (uncompressed binder-bonded
contact, mid-range of literature values for binder-bonded battery
interfaces) and $\alpha = 0.5$~MPa$^{-1}$ (pressure sensitivity, mid of
the 0.3--1.0~MPa$^{-1}$ literature range for compliant
solid-electrolyte/electrode contacts).  These closure parameters
should be re-fit from peel-test and four-point-probe measurements
before quantitative go/no-go decisions are taken.

The contact-stability metric is
\begin{equation}
\Pi_c = \frac{\langle p_c \rangle}{p_{\min}},
\qquad p_{\min} = \frac{1}{\alpha}\,\ln\!\frac{R_0}{R_c^{\rm target}},
\label{eq:Pic}
\end{equation}
with $R_c^{\rm target} = 20$~m$\Omega\cdot$cm$^2$, the cell target.
$\Pi_c \ge 1$ is the design objective: the average contact pressure
exceeds the minimum required to land $R_c$ at or below target.  With
the closure parameters above, $p_{\min} \approx 4.6$~MPa.

\section{Design variables and sweep ranges}
\label{sec:design}

Table~\ref{tab:design-vars} lists the seven design variables and the
range explored in this study.  The first block (eigenstrain magnitudes
and applied pressure) controls the static post-cure compression and is
the focus of the design map below.  The second block (electrode
breathing) controls cycling modulation of $\Pi_c$.  The third block
(geometry and material moduli) is held fixed for the surveyed grid,
with the GPE modulus varied as an additional axis.

\begin{table}[H]
\centering
\caption{Design variables and ranges.}
\label{tab:design-vars}
\begin{tabular}{lccc}
\toprule
Variable & Symbol & Default & Sweep range \\
\midrule
GPE polymerization shrink     & $\varepsilon_{\rm GPE}$  & 0.05 & 0.02--0.12 \\
Package coating shrink        & $\varepsilon_{\rm pkg}$  & 0.05 & 0.02--0.10 \\
External stack pressure       & $\sigma_{\rm stack}$     & 0~MPa & 0--1~MPa \\
\midrule
Anode lithiation expansion    & $\alpha_{\rm Li}^{\rm a}$ & $+0.10$ & --- \\
Cathode lithiation expansion  & $\alpha_{\rm Li}^{\rm c}$ & $-0.05$ & --- \\
Cycling amplitude / period    & $A_{\rm cyc},\,T_{\rm cyc}$ & 0,\,1 & demonstration \\
\midrule
GPE modulus                   & $E_{\rm sep}$            & 10~MPa  & 1--50~MPa \\
Package modulus               & $E_{\rm pkg}$            & 1000~MPa & --- \\
Layer thicknesses             & $t_*$                    & Table~\ref{tab:geom} & --- \\
\bottomrule
\end{tabular}
\end{table}

\section{Results and design-window characterization}

The 192-point sweep covers a four-axis grid
$(\varepsilon_{\rm GPE},\,\varepsilon_{\rm pkg},\,\sigma_{\rm stack},\,
E_{\rm sep})$.  Figure~\ref{fig:heatmap} shows a representative slice
at $E_{\rm sep} = 10$~MPa.  The white $\Pi_c = 1$ contour splits each
panel into a feasible region (right of the contour) and an infeasible
one (left).  Even at $\sigma_{\rm stack} = 0$ (truly pressure-free), a
defined window with $\varepsilon_{\rm pkg} \gtrsim 0.07$ already lands
above target; adding 0.5--1.0~MPa external pressure widens that window
further.  $\Pi_c$ values across the surveyed grid span 0.03--1.94, so
the variables are well-engineered to land near the design target.
Three observations explain the trend structure.

\begin{figure}[H]
\centering
\includegraphics[width=\linewidth]{figures/design_map_heatmap.pdf}
\caption{Contact-stability metric $\Pi_c$ over the
$(\varepsilon_{\rm GPE},\,\varepsilon_{\rm pkg})$ grid at three
externally applied stack pressures (slice at $E_{\rm sep}=10$~MPa).
The white contour marks the $\Pi_c = 1$ design boundary: cells to its
right meet the $R_c \le 20\,\mathrm{m\Omega\,cm^2}$ target with
margin, those to its left fall short.  Even with no external stack
pressure, package shrink $\varepsilon_{\rm pkg}\gtrsim 0.07$ alone
already lands the design above target; modest applied pressure shifts
the boundary to lower $\varepsilon_{\rm pkg}$ values, opening the
feasible window.}
\label{fig:heatmap}
\end{figure}

\subsection{Package shrink dominates compression}

$\Pi_c$ scales nearly linearly with $\varepsilon_{\rm pkg}$ at fixed
$\varepsilon_{\rm GPE}$ and $\sigma_{\rm stack}$
(Fig.~\ref{fig:lines-pkg}); doubling $\varepsilon_{\rm pkg}$
approximately doubles $\langle p_c \rangle$.  The $\varepsilon_{\rm
GPE}$ family of curves is nearly parallel, so $\varepsilon_{\rm pkg}$
acts roughly independently of GPE shrink within this range.  Package
shrink is therefore the primary linear design lever for the
pressure-free regime.

\begin{figure}[H]
\centering
\includegraphics[width=\linewidth]{figures/Pi_c_vs_eps_pkg.pdf}
\caption{$\Pi_c$ vs.\ $\varepsilon_{\rm pkg}$ at four
$\varepsilon_{\rm GPE}$ levels and three external stack pressures
(slice at $E_{\rm sep}=10$~MPa).  Dashed line marks the $\Pi_c=1$
target.  The slope is nearly constant ($\sim$17 per unit
$\varepsilon_{\rm pkg}$); the design crosses target around
$\varepsilon_{\rm pkg}\approx 0.07$ at the pressure-free condition.}
\label{fig:lines-pkg}
\end{figure}

\subsection{GPE shrink can locally unload the separator interfaces}

At fixed $\varepsilon_{\rm pkg}$, increasing $\varepsilon_{\rm GPE}$
\emph{decreases} the contact pressure at the separator--cathode
interface and can drive that interface into tension when
$\varepsilon_{\rm pkg}$ is small (Fig.~\ref{fig:lines-GPE}).
Physically, the GPE shrink pulls the cathode toward the anode through
the soft separator, unloading those interfaces.
Figure~\ref{fig:bars} resolves the effect interface by interface: at
$(\varepsilon_{\rm GPE},\,\varepsilon_{\rm pkg}) = (0.12,\,0.02)$ the
sep$|$cathode interface is in tension---a debonding risk---even though
the stack-averaged $p_c$ remains positive.  The design implication is
that $\varepsilon_{\rm GPE}$ must be chosen \emph{jointly} with
$\varepsilon_{\rm pkg}$, not independently maximized.

\begin{figure}[H]
\centering
\includegraphics[width=\linewidth]{figures/Pi_c_vs_eps_GPE.pdf}
\caption{$\Pi_c$ vs.\ $\varepsilon_{\rm GPE}$ at four
$\varepsilon_{\rm pkg}$ levels.  All curves slope mildly negative:
larger GPE shrink unloads the separator interfaces.}
\label{fig:lines-GPE}
\end{figure}

\begin{figure}[H]
\centering
\includegraphics[width=0.92\linewidth]{figures/interface_pc_bars.pdf}
\caption{Per-interface contact pressure for four representative
design points.  The orange ``low pkg, high GPE'' configuration drives
the sep$|$cathode interface into tension.  Triplets in the legend are
$(\varepsilon_{\rm GPE}/\varepsilon_{\rm pkg}/\sigma_{\rm stack})$.}
\label{fig:bars}
\end{figure}

\subsection{GPE stiffness is a real but diminishing-returns lever}

The fourth-axis sweep over $E_{\rm sep} \in \{1, 10, 25, 50\}$~MPa
(Fig.~\ref{fig:E-sep-lever}) quantifies how stiffening the GPE
amplifies pressure transmission through the stack.  At the
most-favorable shrink corner ($\varepsilon_{\rm GPE}=0.02$,
$\varepsilon_{\rm pkg}=0.10$, $\sigma_{\rm stack}=0$), going from
$E_{\rm sep} = 1 \to 10 \to 25 \to 50$~MPa lifts $\Pi_c$ from
$1.22 \to 1.61 \to 1.81 \to 1.94$.  Returns flatten visibly above
$\sim$25~MPa: once the GPE is stiff enough that the through-stack
modulus is no longer separator-dominated, further stiffening adds
little.  At the soft end ($E_{\rm sep} = 1$~MPa) only the largest
package shrink ($\varepsilon_{\rm pkg} = 0.10$) just clears
$\Pi_c = 1$ when no external pressure is applied; in this regime an
external preload (such as a thermal-shrink overwrap) is the most
effective way to close the gap.  In the practical engineering window,
$E_{\rm sep}$ in the 10--25~MPa range, $\varepsilon_{\rm pkg}$ in the
0.05--0.10 range, and modest $\sigma_{\rm stack}$ together comfortably
hit $\Pi_c \ge 1$.

\begin{figure}[H]
\centering
\includegraphics[width=\linewidth]{figures/Pi_c_vs_E_sep.pdf}
\caption{$\Pi_c$ vs.\ GPE Young's modulus $E_{\rm sep}$ at
$\varepsilon_{\rm GPE} = 0.02$, across $\varepsilon_{\rm pkg} \in
\{0.02, 0.05, 0.08, 0.10\}$ and three external stack pressures.
Diminishing returns set in above $E_{\rm sep} \approx 25$~MPa;
designs with $\varepsilon_{\rm pkg} \ge 0.08$ clear $\Pi_c=1$
across most of the surveyed $E_{\rm sep}$ range.}
\label{fig:E-sep-lever}
\end{figure}

\section{Summary and outlook}

Within the surveyed design grid, three engineering levers control
contact stability of a pressure-free monolithic Li-ion cell:

\begin{enumerate}
\item Package coating shrink $\varepsilon_{\rm pkg}$ is the dominant
linear lever; values in the 0.07--0.10 range alone deliver $\Pi_c \ge
1$ at the pressure-free condition.
\item GPE shrink $\varepsilon_{\rm GPE}$ must be co-designed with
$\varepsilon_{\rm pkg}$.  Independently maximizing GPE shrink is
counter-productive at low package shrink: it locally unloads the
separator--electrode interfaces and creates a debonding risk.
\item GPE modulus $E_{\rm sep}$ provides a useful but diminishing
boost; targets above $\sim$25~MPa offer little additional margin.
\end{enumerate}

Lithiation-induced breathing modulates $\Pi_c$ by $\sim$25\% per
cycle, so the design must clear target at the discharged state---which
is the floor of the cycle for the chosen anode/cathode breathing
parameters.

The headline limitations of the present model are:

\begin{itemize}
\item Mechanics-only.  Adding the charge-balance and Li-conservation
equations brings the second pair of cell-level targets (effective
ionic and electronic conductivity) into the verifiable set, and is
the natural follow-on study.
\item Linear, isotropic homogenization of the porous current
collectors and composite electrodes.  A more faithful model would
distinguish the conductive-network fractions, binder fractions, and
active-material volume fractions.
\item Cohesive-zone behavior at the electrode/separator and
electrode/collector interfaces is not yet modelled.  Once
adhesion-energy estimates from atomistic calculations or peel
measurements are available, debonding under cycling becomes an
additional failure mode to track alongside contact loss.
\item The contact-resistance closure $R_c(p_c)$ uses representative
literature values for $(R_0, \alpha)$.  These should be replaced with
values fit to direct measurements before quantitative go/no-go
decisions are taken from $\Pi_c$.
\item The 2D plane-strain slice with a 1~mm representative width
under-resolves the wrap mechanism by a factor estimated at $\sim$80
relative to the realistic geometry; absolute $\Pi_c$ numbers should
therefore be read as relative indicators along the design axes until
a 3D shell-wrap analysis is performed.
\end{itemize}

\end{document}
